\section{Phase 2 setup: the Ibex CPU and the LOD instruction}\label{sec:ibex}

\subsection{What exactly is simulated}

We did not have a physical chip yet (an FPGA board comes later). Instead, we ran the \emph{real hardware
description} of Ibex in a simulator:
\begin{itemize}
  \item \textbf{Ibex} is written in SystemVerilog, a language that describes circuits: wires, registers,
        adders, and how they are connected.
  \item \textbf{Verilator} translates that description into a C++ program that computes, clock tick by
        clock tick, exactly what every wire of the chip would do. This is called a
        \textbf{cycle-accurate} simulation: the cycle counts are the ones a real chip built from the same
        description would have.
  \item The \textbf{Simple System} is lowRISC's minimal test computer around Ibex: the CPU, 1~MB of
        memory (RAM that answers in one cycle), a timer, and a ``printer'' address. Characters written there
        appear in a log file (\file{ibex\_simple\_system.log}).
\end{itemize}

\begin{kid}
Verilator builds a perfect toy model of the chip inside the computer, so detailed that it moves every
little part once per tick. It is slow (about 1 million ticks per second, while a real chip does 50
million or more), but it never lies about how many ticks a program needs.
\end{kid}

\subsection{The configuration and the tools}

\begin{table}[H]
\centering
\small
\begin{tabularx}{\linewidth}{lX}
\toprule
Ibex version & lowRISC Ibex, fixed at commit \code{7b5df75} (the last one before a big rewrite), plus our
patch \file{hw/patches/0001-ibex-lod-instruction.patch}\\
CPU configuration & the ``small'' configuration, with the multiplier set to \code{RV32MNone}: the core
cannot run \code{mul} or \code{div}. Instruction set RV32IC + our LOD\\
Simulator & Verilator 4.228 (installed with conda, no administrator rights needed), FuseSoC to build\\
Compiler & lowRISC GCC 10.2 for RISC-V (\code{-march=rv32ic -O3}), bare metal (no operating system)\\
Memory & 960~kB for program and data plus a 32~kB stack (\file{hw/sw/link\_1mb.ld}); the int8 MLP weights
alone are 230~kB\\
Timing & the CPU's own cycle counter (\code{mcycle}), read before and after every picture and every
layer\\
Setup scripts & \file{hw/setup.sh} (tools, Ibex, patch), \file{hw/build\_sim.sh} (builds the simulator),
\file{hw/run\_sim.sh} (runs one program)\\
\bottomrule
\end{tabularx}
\caption{The Ibex setup.}
\end{table}

\begin{careful}
The compiler's helper library (libgcc) was built for a core \emph{with} a multiplier, so a few of its
helpers contain \code{mul}/\code{div} instructions, which would crash our core. Every time a program is
built, \file{hw/run\_ibex.py} (and \file{hw/sw/Makefile}) disassemble it and stop with an error if any
\code{mul}, \code{mulh}, \code{div} or \code{rem} instruction is found. The exact multiply in C
(\code{a * b}) becomes a call to libgcc's \code{\_\_mulsi3}, which is a shift-and-add loop: the
paper's baseline.
\end{careful}

\subsection{How LOD was added to Ibex}\label{sec:lodpatch}

A CPU handles every instruction in steps. Ibex has two main stages: \textbf{fetch} (read the instruction
from memory) and \textbf{decode/execute} (work out what the instruction means, read the registers, compute
in the ALU, write the result back). To add an instruction you must teach the \textbf{decoder} to recognise it
and give the \textbf{ALU} (arithmetic logic unit) a circuit that computes it.

\begin{figure}[H]
\centering
\begin{tikzpicture}[>=Stealth,font=\small,
  blk/.style={draw,rounded corners=2pt,align=center,minimum height=1.1cm,minimum width=2.2cm,fill=gray!6},
  new/.style={draw,thick,rounded corners=2pt,align=center,fill=twoit!25,minimum height=0.8cm}]
  \node[blk] (if) {Fetch\\(IF stage)};
  \node[blk,right=0.9cm of if] (dec) {Decoder\\\file{ibex\_decoder.sv}};
  \node[blk,right=0.9cm of dec] (rf) {Register file\\(32 registers)};
  \node[blk,right=0.9cm of rf,minimum width=3.4cm,minimum height=2.4cm] (alu) {};
  \node[anchor=north,font=\small] at (alu.north) {ALU \file{ibex\_alu.sv}};
  \node[new] (lod) at ($(alu.center)+(0,-0.15)$) {LOD: priority\\encoder (new)};
  \node[new,below=0.35cm of dec,font=\scriptsize] (dnew) {opcode 0x0b $\Rightarrow$ ALU\_LOD (new)};
  \draw[->] (if) -- (dec); \draw[->] (dec) -- (rf); \draw[->] (rf) -- (alu);
  \draw[->] (alu.south) -- ++(0,-0.5) -| node[pos=0.25,below] {result written back to \code{rd}} (rf.south);
\end{tikzpicture}
\caption{Where the patch touches Ibex. The orange parts are new; the rest of the core is unchanged.}
\end{figure}

\paragraph{Step 1: choose an encoding.} Every RISC-V instruction is a 32-bit number. The lowest 7 bits
(the \emph{opcode}) say which kind of instruction it is. RISC-V reserves opcode \code{0001011} (hex
\code{0x0b}, called \textbf{custom-0}) for exactly this purpose: vendors may put their own instructions
there, and no standard instruction will ever use it. LOD uses the standard ``R-type'' layout:

\begin{figure}[H]
\centering
\begin{tikzpicture}[font=\small\ttfamily,
  f/.style={draw,minimum height=0.8cm,align=center}]
  \node[f,minimum width=2.6cm,fill=gray!10] (f7) at (0,0) {0000000};
  \node[f,minimum width=1.9cm,fill=gray!10,right=0pt of f7] (rs2) {00000};
  \node[f,minimum width=1.9cm,fill=roba!15,right=0pt of rs2] (rs1) {01011};
  \node[f,minimum width=1.3cm,fill=gray!10,right=0pt of rs1] (f3) {000};
  \node[f,minimum width=1.9cm,fill=roba!15,right=0pt of f3] (rd) {01010};
  \node[f,minimum width=2.6cm,fill=twoit!25,right=0pt of rd] (op) {0001011};
  \foreach \n/\t/\b in {f7/funct7 = 0/31--25, rs2/rs2 = x0/24--20, rs1/rs1 = a1 (x11)/19--15,
                        f3/funct3/14--12, rd/rd = a0 (x10)/11--7, op/custom-0/6--0} {
    \node[above=1pt of \n,font=\scriptsize\rmfamily] {\t};
    \node[below=1pt of \n,font=\scriptsize\rmfamily,gray] {bits \b};
  }
\end{tikzpicture}
\caption{The 32 bits of \code{lod a0, a1}. As one number: $(11\ll15) + (10\ll7) + \texttt{0x0b} =
\texttt{0x0005850B}$.}
\end{figure}

The assembler does not know the name ``lod'', so the C code uses GCC's generic \code{.insn} directive,
which builds an R-type instruction from its fields:
\begin{lstlisting}[language=C]
static inline uint32_t LOD(uint32_t x)
{
    uint32_t r;
    /* R-type, opcode 0x0b (custom-0), funct3 = 0, funct7 = 0, rd = r, rs1 = x, rs2 = x0 */
    asm volatile(".insn r 0x0b, 0, 0, %0, %1, x0" : "=r"(r) : "r"(x));
    return r;
}
\end{lstlisting}

\paragraph{Step 2: name the opcode and the operation} (\file{rtl/ibex\_pkg.sv}). One line adds
\code{OPCODE\_CUSTOM0 = 7'h0b} to the list of opcodes, and one adds \code{ALU\_LOD} to the list of ALU
operations.

\paragraph{Step 3: teach the decoder} (\file{rtl/ibex\_decoder.sv}). When the opcode is custom-0, read
register \code{rs1}, write the register file, and send the operation \code{ALU\_LOD} to the ALU. Any
custom-0 instruction with other \code{funct7}/\code{funct3} values is reported as illegal.
\begin{lstlisting}
OPCODE_CUSTOM0: begin  // RoBA: LOD rd, rs1 (R-type, funct7 = funct3 = 0, rs2 unused)
  rf_ren_a_o   = 1'b1;   // read register rs1
  rf_we        = 1'b1;   // write the result to rd
  illegal_insn = ({instr[31:25], instr[14:12]} == {7'b000_0000, 3'b000}) ? 1'b0 : 1'b1;
end
...
OPCODE_CUSTOM0: begin  // RoBA: LOD rd, rs1
  alu_op_a_mux_sel_o = OP_A_REG_A;
  alu_op_b_mux_sel_o = OP_B_REG_B;
  alu_operator_o     = ALU_LOD;
end
\end{lstlisting}

\paragraph{Step 4: build the circuit in the ALU} (\file{rtl/ibex\_alu.sv}). A priority encoder: go through
the 32 bits from the lowest to the highest and remember the position of the last 1 you saw. This
\code{for} loop is not a program loop that runs over time: it describes 32 small comparisons that all
exist at the same time as wires, so the answer is ready within the same clock cycle.
\begin{lstlisting}
logic [4:0] lod_result;
always_comb begin
  lod_result = '0;                       // LOD(0) = 0
  for (int unsigned i = 0; i < 32; i++)
    if (operand_a_i[i]) lod_result = 5'(i);
end
...
ALU_LOD: result_o = {27'h0, lod_result};  // in the ALU's result multiplexer
\end{lstlisting}

\paragraph{Step 5: test it.} \file{hw/sw/lod\_test.c} runs the new instruction on 0, on every single bit,
on every ``all ones below bit $i$'' pattern, and on 20{,}000 random numbers, and compares each answer with a
slow software loop. All match.

\begin{keyidea}
The whole change to the CPU is about 40 lines: a new opcode name, two decoder cases and a 32-input
priority encoder in the ALU. LOD runs in one cycle, like an addition, and the paper measured that it adds
4.5\% to the power of the multiplierless core.
\end{keyidea}

\subsection{From a trained network to cycles on the chip}\label{sec:flow}

\begin{figure}[H]
\centering
\begin{tikzpicture}[>=Stealth,font=\small,node distance=0.45cm,
  blk/.style={draw,rounded corners=2pt,align=center,minimum height=1.0cm,text width=3.0cm,fill=gray!6},
  py/.style={blk,fill=roba!12}, hw/.style={blk,fill=twoit!18}]
  \node[py] (ck) {trained int8 checkpoint\\(GPU, one per multiplier)};
  \node[py,right=of ck] (exp) {\file{run\_ibex.py}: integer model + images};
  \node[blk,right=of exp] (h) {\file{model.h}, \file{images.h}\\(C headers)};
  \node[hw,below=0.7cm of h] (gcc) {GCC \code{-march=rv32ic -O3}\\+ \code{-D} flags};
  \node[hw,left=of gcc] (elf) {program (ELF), checked for \code{mul}/\code{div}};
  \node[hw,left=of elf] (sim) {Verilator Ibex\\(48 in parallel)};
  \node[blk,below=0.7cm of sim] (log) {log: per picture the answer, logit hash, cycles};
  \node[py,right=of log] (ref) {Python model of the same integer maths};
  \node[py,right=of ref] (csv) {compare $\Rightarrow$ CSV files, figures, slides};
  \draw[->] (ck) -- (exp); \draw[->] (exp) -- (h); \draw[->] (h) -- (gcc); \draw[->] (gcc) -- (elf);
  \draw[->] (elf) -- (sim); \draw[->] (sim) -- (log); \draw[->] (log) -- (ref); \draw[->] (ref) -- (csv);
  \draw[->,dashed] (exp.south) -- ++(0,-0.35) -| ($(ref.north)+(0.6,0)$);
\end{tikzpicture}
\caption{The Phase~2 pipeline. Blue boxes run in Python on the server, orange boxes build and run the
program for the simulated chip.}
\label{fig:flow}
\end{figure}

Step by step, for one dataset, one network and one \emph{configuration} (a multiplier plus options):
\begin{enumerate}
  \item \textbf{Load the checkpoint} fine-tuned at int8 with that multiplier (the table-of-squares runs use the
        exact one).
  \item \textbf{Build the integer model} (\code{int\_model}): int8 weights, $M$ and $C$ per neuron
        (Section~\ref{sec:quant}), written as C arrays into \file{model.h}.
  \item \textbf{Quantize the first 100 test pictures} to int8 and write them, 10 per batch, into
        \file{images.h}.
  \item \textbf{Compute the expected answers in Python}: \code{ref\_forward} repeats exactly the integer
        arithmetic of the C program, bit for bit, and computes a \emph{logit hash} of the 10 scores.
  \item \textbf{Compile} \file{hw/sw/dnn.c} with flags that pick the multiplier and the method, for example
        \code{-DMULT\_NZ=roba\_dn1\_nz} (fused RoBA-Dn1) or \code{-DSQTAB} (table of squares).
  \item \textbf{Simulate}: every batch runs in its own Verilator process, 48 at a time on the server.
        The program prints one line per picture:\\
        {\footnotesize\code{IMG,<index>,<label>,<prediction>,<logit hash>,<total cycles>,<cycles per layer>}}.
  \item \textbf{Compare}: the prediction and the logit hash must equal the Python ones for every picture.
        The summaries go to \file{results/ibex/*.csv}.
\end{enumerate}

\begin{example}[: what a logit hash is for]
The 10 scores of a picture are 64-bit integers, e.g.\ $(-1043, 2210, \dots)$. Comparing all of them for
10{,}400 pictures would make long log files. Instead, the program mixes the 10 numbers into one 32-bit
fingerprint (rotate, XOR, repeat). If even one bit of one score differs, the fingerprint is completely
different. Equal hashes means equal scores.
\end{example}

\paragraph{The four stages of \file{hw/run\_ibex.py}.}
\begin{itemize}
  \item \code{--stage mult}: time one multiplication (10{,}000 of them per routine and operand size).
  \item \code{--stage dnn}: both networks, both datasets, eight multipliers fused into the loop, plus the
        table of squares: 36 configurations.
  \item \code{--stage call}: the same eight multipliers called once per product, as in the paper: 32
        configurations.
  \item \code{--stage zeroskip}: exact, Dn1 and the table, each with/without zero skipping and with
        normalised/raw pixels: 48 configurations.
\end{itemize}
In total: 116 configurations $\times$ 100 pictures = 11{,}600 simulated inferences (10{,}400 different
ones, because 12 zeroskip configurations repeat the dnn ones), about 160 billion simulated clock cycles in
roughly 70 minutes on 48 CPU cores. A finished simulation is reused only if the compiled program is
byte-for-byte the same (its SHA-256 hash matches), so a rerun after an unrelated change is quick.
\newpage
